\begin{scanpage}{0196}{178}
Considering that the quantity $1/D$ is equal to the infinite geometric progression,
\[
\frac1D=1+\psi^2(L)e^{2ikL}+\psi^4(L)e^{4ikL}+\ldots,\tag{28.09}
\]
Expression (28.03) may be written in the expanded form
\[
\begin{aligned}J(z)=J_0\bigl\{&\psi(|z|)e^{ik|z|}-\psi(-z_1)e^{-ikz_1}[\psi(z-z_1)e^{ik(z-z_1)}\\
&-\psi(L)e^{ikL}\psi(z_2-z)e^{ik(z_2-z)}+\psi^2(L)e^{2ikL}\psi(z-z_1)e^{ik(z-z_1)}-\ldots]\\
&-\psi(z_2)e^{ikz_2}[\psi(z_2-z)e^{ik(z_2-z)}-\psi(L)e^{ikL}\psi(z-z_1)e^{ik(z-z_1)}\\
&+\psi^2(L)e^{2ikL}\psi(z_2-z)e^{ik(z_2-z)}-\ldots]\bigr\}.
\end{aligned}\tag{28.10}
\]
The physical meaning of Expression (28.03) is seen from this. The first term in Equation (28.10) is the primary current wave which coincides with the wave excited by a concentrated emf in an infinitely long conductor. The second term (in all brackets) corresponds to the current resulting from the reflection of the primary current wave from the conductor end $z=z_1$, and as a result of subsequent reflections from the conductor ends which arise from this wave $-J_0\psi(-z_1)e^{-ikz_1}\psi(z-z_1)e^{ik(z-z_1)}$. The third term (in all brackets) corresponds to the current resulting from the reflection of the primary wave from the end $z=z_2$ and as a result of the subsequent reflections from the conductor ends arising from this wave $-J_0\psi(z_2)e^{ikz_2}\psi(z_2-z)e^{ik(z_2-z)}$.

It also follows from Equation (28.03) that external field (28.02) excites in the semi-infinite conductor ($z_1<z<\infty$) the current
\[
J(z)=J_0[\psi(|z|)e^{ik|z|}-\psi(-z_1)e^{-ikz_1}\psi(z-z_1)e^{ik(z-z_1)}],\tag{28.11}
\]
and in the semi-infinite conductor ($-\infty<z<z_2$) the current
\[
J(z)=J_0[\psi(|z|)e^{ik|z|}-\psi(z_2)e^{ikz_2}\psi(z_2-z)e^{ik(z_2-z)}].\tag{28.12}
\]
Comparing these expressions with the proper terms in (28.10), we see that the reflection of all the current waves at the end of a finite
\end{scanpage}
